% Generated draft: requires mathematical and source review.
\documentclass[11pt]{article}
\usepackage[margin=25mm]{geometry}
\usepackage{fontspec}
\setmainfont{Latin Modern Roman}
\usepackage{amsmath,amssymb,mathtools,unicode-math}
\setmathfont{Latin Modern Math}
\usepackage{enumitem,needspace,xurl,hyperref}
\setlength{\parindent}{0pt}
\setlength{\parskip}{5pt}
\setlength{\emergencystretch}{3em}
\newcommand{\N}{\mathbb{N}}
\newcommand{\Z}{\mathbb{Z}}
\newcommand{\Q}{\mathbb{Q}}
\newcommand{\R}{\mathbb{R}}
\newcommand{\C}{\mathbb{C}}
\newcommand{\F}{\mathbb{F}}
\newcommand{\A}{\mathbb{A}}
\newcommand{\PP}{\mathbb{P}}
\newcommand{\E}{\mathbb{E}}
\newcommand{\eps}{\varepsilon}
\newcommand{\dd}{\,\mathrm{d}}
\begin{document}
\section*{UnsolvedMath ID: 11000317}
\textbf{Source problem code:} catalogue entry
\subsection{Definitions and mathematical statement}
Retain the source's natural inclusions
\[
 j_A:M_{g,1}\hookrightarrow\operatorname{Aut}(F_{2g}),
 \qquad j_O:M_{g,*}\hookrightarrow\operatorname{Out}(F_{2g}).
\]
The exact boundary, marked-point, or puncture conventions defining these groups and inclusions are absent. In particular, the free group has rank \(2g\); the second map must not be replaced without explanation by the closed-surface fundamental-group action, whose target is a different outer automorphism group.

\emph{Problem 4.11.} Determine whether the induced homomorphisms on reduced rational cohomology are zero:
\[
 j_A^*:\widetilde H^*(\operatorname{Aut}(F_{2g});\Q)
       \longrightarrow\widetilde H^*(M_{g,1};\Q),
\]
\[
 j_O^*:\widetilde H^*(\operatorname{Out}(F_{2g});\Q)
       \longrightarrow\widetilde H^*(M_{g,*};\Q).
\]
These maps are contravariant to the group inclusions. For connected group classifying spaces, reduced cohomology removes the degree-zero unit; triviality asks for vanishing in every positive degree. The two maps and all their intended genera/degrees must be distinguished, and a nonzero class in either source need not have nonzero restriction. The excerpt does not state a genus bound or propose a nonvanishing degree.

The source refers to Wahl [107] for a homotopy-theoretic property of \(j_A\) as \(g\to\infty\), without supplying it. The research summary describes stable mapping-class rational cohomology by Mumford--Morita--Miller classes, crediting Madsen--Weiss, and stable positive-degree \(\operatorname{Aut}/\operatorname{Out}\) rational vanishing, crediting Galatius. These supplied stable results can force reduced source vanishing in a stable range, but do not settle restrictions in every finite genus and unstable degree. No stable theorem is silently extended to all degrees at fixed \(g\).
\subsection{Short English statement}
Are the two restriction maps from rank-\(2g\) free-group automorphism/outer-automorphism cohomology to the specified mapping-class cohomology identically zero in positive degrees? Reduced rational cohomology and finite-genus unstable cases are essential; stable vanishing alone does not answer the full question.
\subsection{Sources}
\begin{itemize}
\item Local UnsolvedMath registry, permanent ID 11000317, source code catalogue entry, Problem 4.11, PDF page 367. No individual proposer is supplied.
\item Farb, \emph{Problems on Mapping Class Groups and Related Topics}. Supplied URL: \url{https://www.math.uchicago.edu/~farb/papers/mcgbook.pdf}; recorded access 2026-07-29; extraction pdf-font-structure. The volume contains 329 explicitly typeset headings across chapters, with repeated labels disambiguated by page.
\item Wahl [107] is cited for a stable homotopy-theoretic property of the first inclusion. Its full entry and theorem are absent. The summary credits Madsen--Weiss and Galatius with stable results but gives no full bibliography or stability bounds.
\item Supplied status open and OPEN-TRIAGE. The summary says the finite-genus reduced-cohomology question appears unresolved. Current primary-literature status and redistribution terms remain NEEDS\_REVIEW.
\end{itemize}
\end{document}
